L'espace est muni d'un repère orthonormé \Oijk{} d'unité 1 cm.

\subsection*{Partie A}

On considère les points de l'espace R$(1 ; 0 ; 6)$, S$(1 ; 3 ; 2)$, T$(\np{3,5} ; 7 ; 5)$ et A$(\np{-2,5} ; 8 ; \np{4,5})$.

\begin{enumerate}
	\item Démontrer que les points R, S et T ne sont pas alignés.

	\item \begin{enumerate}
		\item Démontrer que le vecteur $\vect{n} \begin{pmatrix}
			-10 \\ 4 \\ 3	\end{pmatrix}$ est un vecteur normal au plan (RST).

		\item En déduire une équation cartésienne du plan (RST).
	\end{enumerate}
	\item \begin{enumerate}
		\item Vérifier que la point A n'appartient pas au plan (RST).

		\item On note $d_1$ la droite passant par le point A et orthogonale au plan (RST).

		Déterminer une représentation paramétrique de la droite $d_1$.

		\item Soit le point H projeté orthogonal de A sur (RST).

		Vérifier que ses coordonnées sont $(\np{2,5} ; 6 ; 3)$.
	\end{enumerate}

	\item \begin{enumerate}
		\item Démontrer que le triangle RST est rectangle en S.

		\item Vérifier que le volume du tétraèdre ARST est égal à $\dfrac{625}{24}$ cm\up{3}.

		On rappelle que le volume du tétraèdre est donné par la formule :
		\[ V= \frac{1}{3} \times \text{aire de la base} \times \text{hauteur}\]
	\end{enumerate}
\end{enumerate}

\subsection*{Partie B}


Soit M$(x ; y ; z)$ un point de la droite $d_2$ de représentation paramétrique :

$$\begin{cases}x=\np{-4,5} +s \\
	y=4-4s \\
	z=\np{5,5}-2s \end{cases},\: s \in \R$$

\begin{enumerate}
	\item Démontrer que $\mathrm{AM}^2 = 21s^2 + 24s + 21$.

	\item Déterminer les coordonnées du point M de la droite $d_2$ pour lesquelles la distance AM est minimale.
\end{enumerate}
